% Generated by task tex-changes from dossiers/register.md (rows R1 to R140) and the finding boxes
% of gen/findings-*.tex. Every row points to the text that governs it; nothing here is new.
% No \chapter line: a-proposed-changes.tex provides it. Blueprint lines are at b435631
% (at 144c5aa every line after 51 is four higher).
\label{sec:chg}

\paragraph{How to read this appendix.} Each row of \cref{tab:chg-blueprint} is one change to the
blueprint (\file{docs/roadmap/protocol-blueprint.md}), placed where it applies: the rows follow the
blueprint's own order, section by section (the headline and \enquote{For zkVM engineers}, the scope,
the contracts and the ledger, the pinned conventions row by row, the spine, Layers 0 to 13, the
acceptance tests, the interfaces, the boundaries, the ordering, how work is tracked, the
references), and a finding that touches passages in several sections has a row in each. The first
column names the passage by its section and its lines at leanerVM \code{b435631} (at \code{144c5aa}
every line after 51 is four higher; the lines were re-read at \code{b435631} for this table, and
where a dossier's line differs the table gives the right one). The second column says the change
in a sentence or two. It is a summary: the text that governs is the one the third column points
to, the register row (\code{R1} to \code{R140}, \cref{sec:reg}) and the finding box that gives the
passage as it stands, as proposed and the reason, or, where no box has been typeset, the section
of the dossier under \file{.claude/reports/blueprint-review/dossiers/}. The fourth column is the
classification the dossier gives, in one of four words: an error of the blueprint, a deliberate
deviation, a deviation forced by an upstream library, or documentation (stale, missing or
misplaced text that is no divergence from leanVM); where the dossier gives none, the column says
so. The last column is the severity of the strongest finding behind the row.

Three kinds of change are not in \cref{tab:chg-blueprint}. Changes to the status, the tracker, the
hole comment, \file{docs/architecture.md}, \file{docs/leanvm-target.md} and the other documents
are in \cref{ch:documentation} and in the drafts of \cref{sec:docs-texts} (among them the status
rows \code{R38}, \code{R83}, \code{R94}, \code{R95}, \code{R97} and the attribution of
\code{R100}). Findings against leanVM's own sources, to be reported to leanVM or recorded in
\file{docs/leanvm-target.md} (\code{R105} to \code{R109}, \code{R111}, \code{R112}, \code{R116} to
\code{R118}), change nothing here. Changes to the code are in the shorter
\cref{tab:chg-code}, with the file and lines at \code{b435631} in place of the passage. The rows
that ask the owner to choose between alternatives are marked \textsf{[decision]}; the alternatives
and the review's recommendation are in \cref{sec:chg-decisions}. The register's positions \code{R3}
and \code{R8} and its rows that propose no change to the blueprint (\code{R101}, \code{R104},
\code{R110}, \code{R113}, \code{R115}, \code{R123}, \code{R126}, \code{R133}, \code{R134},
\code{R136}, \code{R140}) appear only where they bear on a row.

\section{Changes to the blueprint}\label{sec:chg-blueprint}

\begin{landscape}
\footnotesize
\setlength{\tabcolsep}{4pt}
\renewcommand{\arraystretch}{1.12}
\begin{longtable}{@{}>{\raggedright\arraybackslash}p{3.2cm}>{\raggedright\arraybackslash}p{10.2cm}>{\raggedright\arraybackslash}p{4.6cm}>{\raggedright\arraybackslash}p{3.1cm}>{\raggedright\arraybackslash}p{1.6cm}@{}}
\caption{Proposed changes to the blueprint, in the blueprint's order (lines at \code{b435631}).}\label{tab:chg-blueprint}\\
\toprule
\textsf{\bfseries Passage} & \textsf{\bfseries Change} & \textsf{\bfseries Register row; full text} & \textsf{\bfseries Classification} & \textsf{\bfseries Severity}\\
\midrule
\endfirsthead
\multicolumn{5}{@{}l}{\textit{\Cref{tab:chg-blueprint}, continued.}}\\
\toprule
\textsf{\bfseries Passage} & \textsf{\bfseries Change} & \textsf{\bfseries Register row; full text} & \textsf{\bfseries Classification} & \textsf{\bfseries Severity}\\
\midrule
\endhead
\midrule
\multicolumn{5}{r@{}}{\textit{Continued on the next page.}}\\
\endfoot
\bottomrule
\endlastfoot
%
\multicolumn{5}{@{}l}{\textsf{\bfseries Throughout the blueprint}}\\*
Throughout (codes in \code{:260}, \code{:592-593}, \code{:1440} and passim) &
Holes, ledger entries and findings are named in words; layers, acceptance tests, decisions and
target theorems keep their numbers with the name beside them; one index of names at the end of the
blueprint. The two references that resolve to the wrong finding are rows below. &
\code{R85}; \cref{f:docs-letter-codes}; \cref{sec:docs-proposal-codes} &
documentation & \sev{minor}\\
Throughout (the ledger, the holes' specifications, the bus phase's four statements) &
One home per kind of fact: the blueprint holds what is wanted and decided (targets, conventions,
the specification of every hole, the acceptance tests, the decisions, the upstream ledger); the
status, the tracker and the hole comment link to it and restate nothing. &
\code{R86}; \cref{f:docs-many-homes}; \cref{sec:docs-proposal-owners} &
documentation & \sev{minor}\\
\midrule
\multicolumn{5}{@{}l}{\textsf{\bfseries The headline and \enquote{For zkVM engineers} (\code{:1-103})}}\\*
Headline, \code{:8-17} &
The two pseudo-signatures carry the bundle of phases and its proofs (\lean{P}, \lean{C},
\lean{S}), the named extractor \lean{piopExtractor P S} and the error \lean{piopError I}; a
sentence says both are composition theorems, which hold of every bundle meeting the seams and say
nothing of leanVM until the instance and the five phases are leanVM's, and that \code{main} proves
the spine, the commit phase and the public-input phase. &
\code{R45}; \cref{f:code-spine-headline-describes-theorems-that-do} &
error of the blueprint & \sev{minor}\\
Headline \code{:21}; For zkVM engineers \code{:77} &
\enquote{the polynomial view of any Clean ensemble} and \enquote{which \lean{Ensemble.toM3}
derives from any Clean \lean{Ensemble}} go (or \lean{Ensemble.toM3} is defined as the six-table
part only): the leanISA instance is not \lean{Ensemble.toM3} of the eight tables, and no map from
an \lean{Ensemble} to an \lean{M3Instance} exists (an instance needs $μ$, a layout, public lines,
\lean{aux}, \lean{Stmt}). &
\code{R30}; \file{boundary-adaptor.md} §G.1 &
error of the blueprint & \sev{major}\\
For zkVM engineers, \code{:62} &
Round-by-round \emph{knowledge} soundness is from [CMS19] Definition 8.5, not \enquote{the standard
form since [CCHLRR19]} (the reference entry itself is corrected under References). &
\code{R98}; \file{literature.md} §A.4 (fifth finding) &
error of the blueprint & \sev{minor}\\
For zkVM engineers, a paragraph after \code{:94} &
State the scope of the random-oracle assumption: the non-interactive theorems are in the
random-oracle model for the one 64-byte BLAKE2s map the chain, the Merkle trees and the grinding
share; they do not cover attacks that use the code of BLAKE2s ([KRS25], [Fen26]); leanVM binds
the statement and commits the trace before the first challenge, so the known attacks do not
apply, and no theorem says more. &
\code{R67}; \file{literature.md} §B.4 &
error of the blueprint (omission) & \sev{minor}\\
For zkVM engineers, \code{:99-101} &
\enquote{reads the stack \lean{q} off the oracle message} is true of the commit phase's
intermediate extractor at round 0, not of the composed extractor's output (the last phase's);
the extracted stack becomes a definition, \lean{piopExtractedStack}, to which the adaptor applies
\lean{witnessOf}. &
\code{R46}; \cref{f:code-spine-extracted-stack-not-definition-rfl} &
error of the blueprint & \sev{minor}\\
\midrule
\multicolumn{5}{@{}l}{\textsf{\bfseries Scope (\code{:105-158})}}\\*
In scope, \code{:109} (\code{:113} at \code{144c5aa}) &
\enquote{ArkLib as a Lake dependency at \code{dca90385}} becomes \enquote{ArkLib as a Lake
dependency, with \dots}: the pin moved at \code{144c5aa}, and pins are stated only in
\file{upstreams.json} and \file{docs/dependencies.md}. &
\code{R84}; \cref{f:docs-superseded-pins} &
documentation (stale at \code{144c5aa}) & \sev{minor}\\
In scope, \code{:126-128} &
\enquote{the claim pool and the opening sumcheck} loses \enquote{the opening sumcheck}: with the
inner-product oracle the opening phase is the batching challenge and one weighted query
(specification §8.5, \enquote{there is no separate reduction sumcheck}). \textsf{[decision]} &
\code{R18}; \file{gt-opening-compile.md} §I.1, §A.5 &
error of the blueprint (caused by Layer 0's evaluation interface) & \sev{major}\\
Out of scope, \code{:153-155} &
Say that [BCHKS25] Theorem 4.6, the one coding-theory input, is a preprint theorem whose proof is
sketched (§4.3; the full argument for the quadratic bound in [Hab25]), peer review pending as of
2026-09. &
\code{R114}; \file{literature.md} §C.4 (first finding) &
documentation (a fact about the assumption) & \sev{note}\\
Out of scope, after \code{:157} &
Add: security against quantum adversaries (the quantum random-oracle model) is out of scope; the
classical bound does not transfer (the loss is quadratic in the query count). &
\code{R68}; \file{literature.md} §G.5 (second finding) &
error of the blueprint (omission) & \sev{minor}\\
Out of scope, a new item &
Record that adaptive statement soundness in the random-oracle model (the input seeds the chain)
is the deployment theorem's (T7) to state, and is not given by this roadmap's per-statement
theorems, whose games quantify over the input statement outside the probability. &
\code{R103}; \file{boundary-adaptor.md} §G.17 &
none given & \sev{note}\\
